\documentclass[11pt]{article}
\usepackage[T1]{fontenc}\usepackage{lmodern,amsmath,amssymb,amsthm,booktabs,longtable}
\usepackage[margin=.9in]{geometry}\usepackage[hidelinks]{hyperref}
\newtheorem{theorem}{Theorem}\newtheorem{lemma}{Lemma}
\newcommand{\dd}{\,d}\newcommand{\E}{\mathbb E}
\title{Original source channels, producers and signed lag}
\author{Independent analytical increment}\date{Version 1, 2 November 2026}
\begin{document}\maketitle
\begin{abstract}
On the initialized deep-regime learned windows, the six core motion
channels and their fixed-label producer contributions converge at order
\(q^{1/20}\) after normalization by \(q^2\). The exact full-output lag is
positive, and cluster-1 help equals exposure times that signed lag with
a smaller remainder. The weak-produced cluster-2 input sign is established
on \(J\cap(-\infty,0]\), but not on the remainder of J. The increment stops
at this sign-extension obligation; it does not claim a complete B-core
sign theorem on all of J.
\end{abstract}
\section{Scope and inventory}\label{an01cpdv1:scope}
Use the original \(S_0=q^{10}\) family, RTA's unchanged rectangle
\(\mathcal P_*\subset[13/20,7/10]\times[87/50,7/4]\), and its fixed
window \(J=[s_*-r,s_*+r]\). Keep the clock
\[
\tau=\mu_1^2t_{\rm phys},\quad
\tau_h=t_n+\lambda^{-1}\log((1-n)/n),\qquad s=\tau-\tau_h,\quad\lambda=\rho^2.
\]
All statements are uniform on this rectangle/window; q is chosen last.
Write \(d\mu=m\,d\lambda_0\), with FWC's fixed core \(K\), selected
weak cohort S, other weak cohort \(P_q\cup\Sigma\), and outside B.
Comparison coordinates are denoted \(X,Y,U,V\), weak mass \(n_\rho=n\)
in formulas below, and \(A_*=\zeta-X\), \(Z=X-Y\).
The entry level in the clock formula is distinct from this time-dependent n.
Let R bound the fixed strong and comparison charts, and put \(\chi=1/20\).

The exact identities used from P are motion, projection, factorizations,
epsilonB, helpdefect and producer. Their hypotheses are finite Gaussian
moments, balanced characteristic evolution, and the probe chain rule
(or selector/a.e. interpretation). The core factorization additionally
requires positive core mass and that the core be currently probe-active;
RTA coregap gives the strict margin \(g_*/8\), and FWC gives \(M\ge1/2\).
The Gaussian wrong-gate estimate in helpdefect requires positive first
input coordinate, supplied by the compact core chart. No identity asserts
that a remainder is small.

P's matrices \(R_p=\E_{P_p}[(X_{\rm data}-f)X_{\rm data}^T1_{\rm gate}]\)
are unhalved. RTA uses \(R^{(p)}=R_p/2\).
Thus \(v_{\theta,p}=b^TR_pa^\perp/2\),
\(v_{\psi,p}=(b^\perp)^TR_pa/2\), and \(h_p=2g_p=b^TR_pa\).
All output derivatives below use this full radial reaction.

P leading, pop-sources and atomicrates are leading-system statements,
not finite-noise reductions. They require the defined characteristic
system and sufficient moments; atomicrates further uses the selected
coherent branch, \(M=1\), \(0<n<1\), active strong and inactive weak
probe. Their original-flow uses below are supplied by FWC and direct
Gaussian calculations, not by formal substitution.
P lagtheorem requires bounded leading charts, \(M=1\), \(0\le N\le1\),
positive incoming lag and finite coordinates for strict Xi. It is not
used to infer exact lag positivity. P tracking requires fixed initial
labels and projected kernel evolution; its damping/forcing comparison
has additional quantitative hypotheses. That route is not used.

FWC main, charts, charges, corefield, massfine, distbound and background
are unconditional for this initialization on each fixed finite continuation
window with bounded offsets. Choose such a window covering J.
RTA coregap, corerate and backgroundrate apply on exactly
\(\mathcal P_*\times J\). WIT uniformtheorem supplies its margins
and comparator gap; regiontheorem is used only for \(s\le0\).
P limitreversal's forward invariants \(V>0,V>Y,V>X\) hold on the
selected branch for \(0<\rho<1\), at every finite time; their boundary
proof is recalled in Section 5. No compact weak support is assumed.

\section{Exact channels and producer limits}\label{an01cpdv1:channels}
At the actual state set
\[
z=a\cdot\xi_q,\quad v=-a^\perp\cdot\xi_q,\quad
c=e_\xi\cdot b,\quad d=-e_\xi\cdot b^\perp,\quad e_\xi=\xi_q-f(\xi_q).
\]
Here c is a probe projection, not the weak response. Define
\[
I_p=\int_K vc\,v_{\theta,p}\dd\mu,\quad
O_p=\int_K zd\,v_{\psi,p}\dd\mu,\quad
P_p=-\int_K h_pzc\dd\mu.
\]
P motion gives exactly \(D_p^\tau=I_p+O_p+P_p+D_{p,K^c}\).
For each training producer \(j\), put
\(e^{\rm tar}(X_{\rm data})=X_{\rm data}\) and \(e^j=-f_j\).
Use these in \(R_p\) to define \(I_p^j,O_p^j,P_p^j\).
Then each channel is the sum of these contributions. In this identity
the actual probe residual, receiver state, measure and gate are held fixed.
It is not an assertion that the rate is linear in the full network state.
This extends P producer without changing its rate weighting.

\begin{theorem}[Channel and producer table]\label{an01cpdv1:table}
Each entry in the following table is the limit of the corresponding
exact channel divided by \(q^2\), with uniform error \(O(q^\chi)\).
Set \(\Phi_x=\Phi(\rho X)\), \(K_x=K(\rho X)\), \(K_u=K(U)\).
Each omitted-producer row \(P_q\cup\Sigma\) and B has zero limit.
\[
\begin{array}{c|ccc}
 &I_1/q^2&O_1/q^2&P_1/q^2\\ \hline
{\rm target}&-ZX/2&-A_*Y/2&-A_*Z\\
K&ZX/2&0&A_*Z\\
S&0&-A_*nK_u/2&0\\ \hline
{\rm sum}&0&-A_*(Y+nK_u)/2&0
\end{array}
\]
\[
\begin{array}{c|ccc}
 &I_2/q^2&O_2/q^2&P_2/q^2\\ \hline
{\rm target}&Z[\rho\phi(\rho X)+\lambda Y\Phi_x]/2&
 A_*\rho K_x/2&0\\
K&-Z\rho K_x/2&0&0\\
S&Z\lambda n(V-Y)\Phi_x/2&-A_*\rho nK_x/2&0
\end{array}
\]
In addition \(|D_{p,K^c}|\le Cq^2\epsilon_{\rm bg}\), where
\(\epsilon_{\rm bg}=q^8+(1+\log(1/q))^2q^{289169/924000}\).
In particular \(I_1,P_1=O(q^{2+\chi})\), although their separate
target/core producer contributions are of order \(q^2\).
\end{theorem}
\begin{proof}
BC's background calculation applies source by source: FWC gives
\(|f_{\Sigma\cup B}(\xi_q)|\le Cq\epsilon_{\rm bg}\), strict primary
inactivity, and RTA's full-rate versus core-self error.
The total core output derivative is \(O(q)\). Hence the full versus
self core-rate difference is \(O(q^2\epsilon_{\rm bg})\), and subtraction
also bounds the actual non-core rate asserted here.

First consider the channel sums. Write \(\mathcal A=\int_K(\zeta-x)d\mu\).
On the reached core \(vc/q=\zeta-\mathcal A-y+O(q^2+\epsilon_{\rm bg})\),
\(zd/q=\zeta-x+O(q^2+q\epsilon_{\rm bg})\), and
\(zc/q^2=(\zeta-x)(\zeta-\mathcal A-y)+O(q^2+\epsilon_{\rm bg})\).
These follow from RTA's core sine expansion and the background output bound.
FWC gives \(x'_1=-(1-M)x/2+O(q^{1/10})\),
\(h_1=1-M+O(q^2+\mu_B)\), \(1-M=O(q^{11/10})\), and
\(2y'_1=-(1-M)y-\int_Ky\,d\mu-\int_SK(u)\,d\mu+O(q^{1/10})\).
Multiplying by the three weights and integrating proves the source-1 sum
limits by distbound. For source 2 use BC twoway and its separate output
expansion, with \(h_2=O(q^2)\). These give respectively
\(Zx'_{2,*},\,A_*\rho(1-n)K_x/2,\,0\) for the three limits.

For producer resolution we use an explicit donor comparison that avoids
any compact-support requirement on the weak tail. For one unit feature
\(b(a\cdot X_{\rm data})_+\), changing the scaled coordinates by
\((\Delta x,\Delta y)\) changes the feature pointwise in norm by at most
\(q(|\Delta x|+|\Delta y|)|X_{\rm data}|\), also in the weak chart with
its reversed output convention. Integrating over either fixed donor
cohort and applying FWC distbound gives an \(L^2(P_p)\) error
\(Cq^{1+\chi}\) on replacing its output by a single feature at its
comparison coordinates with the \emph{actual} mass \(M\) or \(N\).
For each actual core receiver, Cauchy--Schwarz against the gated data
therefore costs \(O(q^{1+\chi})\) in every unscaled source matrix,
hence \(O(q^\chi)\) even in scaled angular velocities. This estimate
compares donor outputs inside the integral. The receiver and its
probe weight have not moved.

Here are the explicit Gaussian calculations for these two comparison
donors, still at an actual receiver \((x,y)\). Entries are
\((2x'_p,2y'_p,h_p)\), with errors \(O(q^\chi)\):
\[
\begin{array}{c|ccc}
 &2x'_1&2y'_1&h_1\\ \hline
{\rm target}&-x&-y&1\\
K&Mx&M(y-Y)&-M\\
S&0&-NK(U)&0
\end{array}
\qquad
\begin{array}{c|ccc}
 &2x'_2&2y'_2&h_2\\ \hline
{\rm target}&\rho\phi(t)+\lambda y\Phi(t)&\rho K(t)&0\\
K&-M\rho F(t,\rho X)&0&0\\
S&\lambda N(V-y)\Phi(t)&-N\rho K(t)&0
\end{array}
\]
Here \(t=\rho x\). On \(P_1\) write the data as \((1+qZ_1,qZ_2)\).
The core receiver and comparison core donor gates can be opened with
exponential error by Cauchy--Schwarz and their Gaussian tails. The core
donor output is \((MX_1,qMYX_1)+O_{L^2}(q^2)\), and the weak output is
\((0,qN(Z_2+U)_+)+O_{L^2}(q^2)\), by the Lipschitz positive part.
Contracting against \(a^\perp=(-qx,1)+O(q^2)\),
\(b^\perp=(-qy,1)+O(q^2)\) gives the first table, using
\(\E Z_2=0\), \(\E(Z_2+U)_+=K(U)\).
On \(P_2\) write the data as \((qZ_1,\rho+qZ_2)\).
The core donor output is \((qM(Z_1+\rho X)_+,0)+O_{L^2}(q^2)\);
the weak output is \((-qN\rho V,NX_2)+O_{L^2}(q^2)\).
Its own gate has exponentially small closed tail. For the input
velocity contract \(b\cdot e^j\) against \(a^\perp\cdot X_{\rm data}\)
on the actual receiver gate, dividing by q. For the output velocity
contract \(b^\perp\cdot e^j\) against \((a\cdot X_{\rm data})_+/q\).
The latter feature differs from \((Z_1+t)_+\) by \(O_{L^2}(q)\).
For the former, conditional on \(Z_2\), the gate threshold in \(Z_1\)
is \(-(\tan(qx)/q)(\rho+qZ_2)\), differing from \(-t\) by
\(O_R(q|Z_2|+q^2)\). Integrating over the intervening strip costs
\(O_R(q)\) in each leading Gaussian moment: its integrand is bounded
by a polynomial in \(|Z_2|\) times a bounded polynomial-weighted
normal density in \(Z_1\). Thus the target input uses
\(\E[(Z_1+\rho y)1_{Z_1>-t}]=\phi(t)+\rho y\Phi(t)\);
the core uses \(F(t,\rho X)=\E[(Z_1+\rho X)_+1_{Z_1>-t}]\).
The weak first output gives \(+\lambda NV\Phi(t)\), its second output
gives \(-\lambda Ny\Phi(t)\). Every radial source-2 entry vanishes
in the limit since \(a\cdot X_{\rm data}=O_{L^2}(q)\).

For omitted \(A=P_q,\Sigma\), FWC's energy bound on the crossing
component gives \(q^{-1}\|f_{A,2}\|_{P_1}+
q^{-1}\|f_{A,1}\|_{P_2}\le C(\mu_A+q^{-1}\sqrt{\mu_AJ_A})\)
by its outputenergy inequalities. The other component has norm
\(O(\mu_A)\). On \(P_1\) it enters a scaled angular contraction
with either \(b_1^\perp=O(q)\) or
\(\|a^\perp\cdot X_{\rm data}\|_2=O(q)\); on \(P_2\) use
\(b_2=O(q)\) or \(\|a\cdot X_{\rm data}\|_2=O(q)\).
The radial charge is \(O(\mu_A)\). Hence each omitted row costs
\(O(q^{1/8}+q^{\gamma_*/2})\), \(\gamma_*=289169/924000\).
B costs \(C\mu_B/q=O(q^{1/10})\) by the global feature bound.

Finally let \(\mathcal A=\int_K(\zeta-x)d\mu\). On the active core,
\[
z=q(\zeta-x)+O(q^3),\quad v=1+O(q^2),\quad
c=q(\zeta-\mathcal A-y)+O(q^3+q\epsilon_{\rm bg}),\quad
d=1+O(q^2+q\epsilon_{\rm bg}).
\]
Thus the three normalized channel integrands are respectively
\((\zeta-\mathcal A-y)x'_p\),
\((\zeta-x)y'_p\), and
\(-h_p(\zeta-x)(\zeta-\mathcal A-y)\), up to vanishing errors.
Integrate first against the actual core measure. FWC distance and
RTA supremum tracking then replace these bounded Lipschitz functions
and donor moments by their coherent values at error \(O(q^\chi)\).
The weak function K is globally Lipschitz, so weak first moments suffice.
This proves the tables without evaluating a rate at a mean angle.
\end{proof}

\section{Signed exact lag and help}\label{an01cpdv1:lag}
\begin{theorem}\label{an01cpdv1:help}
For the exact full output put
\(\mathcal L_q=r_{21}/q\), \(Q_1=1+q^2\). Then
\[
\mathcal L_q=Y+nK(U)+O(q^\chi)=L_0+O(q^\chi),\qquad
L_0\ge c_L:=\frac{5p_*}{4(7/4+R)}>0.
\]
With the exact scaled-angle exposure
\(\mathcal A_q=\int_K(\zeta-x)d\mu\),
\[
q^{-2}D_1^\tau=-\tfrac12\mathcal A_q\mathcal L_q+O(q^\chi).
\]
The remainder is smaller than half the positive margin
\(\mathcal A_qc_L/2\) for sufficiently small q.
P's exact \(B_0\) obeys
\(q^{-2}B_0=\mathcal A_q\mathcal L_q/2+O(q^2+q\epsilon_{\rm bg})\).
\end{theorem}
\begin{proof}
Set \(\mathcal Y=\int_Ky\,d\mu,\ K_w=\int_SK(u)\,d\mu\) and
\(p=603/640\) (FWC's weak support exponent).
The target cross moment vanishes exactly. For the core, open its
\(P_1\) gate at exponential cost. The surviving regression coefficient
is \(\int_K\sin(qy)\cos(qx)d\mu\), up to that tail.
For S put \(w=\sin(qu)/q\). Gaussian projection gives its exact
coefficient before division by \(Q_1\) and after division by q as
\(\int_S\cos(qv)[K(w)+q\sin(qu)\Phi(w)]d\mu\).
The 1-Lipschitz property of K and Taylor bounds charge the difference
from \(K_w\) by \(Cq^2\int_S(1+|u|+|v|)^3d\mu=O(q^{2-p})\)
using FWC charts (bounded second moment and supremum \(O(q^{-p})\)).
Consequently
\[
\frac{\E_{P_1}[X_1 f_{K,2}]}{qQ_1}=\mathcal Y+O(q^2),\qquad
\frac{\E_{P_1}[X_1 f_{S,2}]}{qQ_1}=K_w+O(q^{2-p}+q^2).
\]
Other weak and outside terms have the same energy charges as above.
FWC distance therefore proves the signed comparison, including all labels.
This is a direct moment proof, not an absolute tracking estimate.

WIT gives \(-A_*L_0/2=d_{1,*}\le-5p_*/8\).
Since \(0<A_*\le7/4+R\), the stated lower bound follows.
The sharper proposed \(5p_*/7\) would additionally require \(X\ge0\)
on all of J; that fact is not assumed here.
RTA's gap and \(M\ge1/2\) give \(\mathcal A_q\ge g_*/16\).
The source-1 channel table, the lag comparison and exposure comparison
prove the help law. Choose the fixed error constant times \(q^\chi\)
below \(g_*c_L/64\), which is half the available margin.

For the B0 comparison,
\(M\bar z=q\mathcal A_q+O(q^3)\), \(\bar d=1+O(q^2+q\epsilon_{\rm bg})\) by the full probe bound. Thus
\[
q^{-2}B_0=\tfrac12\mathcal A_q\mathcal L_q
 +O(q^2+q\epsilon_{\rm bg}).
\]
Here \(q\epsilon_{\rm bg}\) need not be \(O(q^2)\); retain it.
Accordingly the B0 agreement asserted and used is at \(O(q^\chi)\)
precision. The \(1-M\) input/radial reactions are in the remainder,
not the leading help.
\end{proof}

\section{Signs and the stopping point}\label{an01cpdv1:signs}
Combining target and core input contributions, define
\[
T_{\rm in}=-\tfrac\lambda2(X-Y)^2\Phi(\rho X),\quad
W_{\rm in}=\tfrac\lambda2 n(X-Y)(V-Y)\Phi(\rho X),
\]
\[
T_{\rm out}=\tfrac\rho2 A_*K(\rho X),\qquad
W_{\rm out}=-\tfrac\rho2 nA_*K(\rho X).
\]
The reviewer's four-part table is correct, with target/core combined
only in \(T_{\rm in}\). Uniformly,
\[
q^{-2}D_2^\tau=T_{\rm in}+W_{\rm in}+T_{\rm out}+W_{\rm out}
 +O(q^\chi).
\]
Thus \(T_{\rm in}\le0,T_{\rm out}>0,W_{\rm out}<0\) on all of J.
The output sum is BC's strictly positive \(\mathcal O_*\).
P limitreversal proves \(V>0,V>Y,V>X\) forward: at \(V=0\),
\(2V'=\rho K(\rho X)>0\); at \(V=Y>0\),
\[
2(V-Y)'=\rho nK(\rho X)+(1-\lambda)Y+nK(U)>0;
\]
at \(V=X\), using the previous invariant,
\[
2(V-X)'=\lambda[K(\rho X)/\rho-X+
                       (1-n)(X-Y)\Phi(\rho X)]>0.
\]
These signs start strictly at the resident end and hold for every finite
time by P's forward existence. They do not prove \(X>Y\).
WIT regiontheorem does give \(X>Y\) for \(s\le0\), so
\(W_{\rm in}>0\) on \(J\cap(-\infty,0]\).

Here the direct invariant extension is incomplete. Set \(z=X-Y\),
\(T=Y+nK(U)-\rho(1-n)K(\rho X)\). At the boundary \(z=0\),
\[
2z'=\lambda n(V-Y)\Phi(\rho Y)+T.
\]
The exact additional boundary inequality needed by this route is
\begin{equation}\label{an01cpdv1:missing}
T>-\lambda n(V-Y)\Phi(\rho Y)\quad\text{at every possible }
z=0\text{ boundary after }s=0\text{ up to }\sup J.
\end{equation}
WIT's positive-T estimate assumes \(n\le1/2\) and its half-learning
region. It cannot be applied after that boundary. No bound in the
present proof supplies \eqref{an01cpdv1:missing}.
We therefore stop at Step 6: the sign of \(W_{\rm in}\) on
\(J\cap(0,\infty)\) is unproved here. No missing sign is inserted as
a hypothesis of a full theorem. WIT's construction by the first change of sign before
the probe gate hit fixes \(s_*>s_a\) but supplies no
\(s_*+r\le0\) bound. Thus window placement alone does not discharge
this inequality. On every compact part where a limiting sign is strict,
continuity gives a positive absolute minimum and the uniform channel
error transfers that sign for small q. In particular both output signs
transfer on all of J. Near a potential zero of \(X-Y\), the
nonpositive limiting square alone gives no strict finite-q sign.

The closed mechanism statement consists of the channel tables, positive
exact lag and help law, and the four-part second-source identity with the
signs just stated. RTA still gives the total second-source positivity
and the ordered total-rate windows. The present decomposition does not
single out one evolving factor as their cause.

\section{Constant order and handoff}\label{an01cpdv1:handoff}
Keep WPT/RB--WCS and RTA/FWC constants, WIT's rectangle and J, and a
fixed continuation radius R first. Fix the Gaussian/moment error constants,
\(c_L\) and the channel error constants next; choose q last for all
source thresholds and the displayed help margin. Physical rates multiply
all normalized rates by \(\mu_1^2\). No numerical premise is used.

Cluster 2 has both input and
output rotation of the strong carrier. Selected weak output changes both
the first-coordinate input drive and the second-coordinate response;
its input term is \(W_{\rm in}\), its output term \(W_{\rm out}\).
Target/core input opposes, and target output supports, the second-source
rate. A pure input-rotation description is incomplete.

The unresolved extension is the sign condition above, not the exact
producer identity. Canonical/bridge coverage still requires initialized
access to that regime and remainder precision relative to its small
exact lag cancellation. No H2, P, README or prior increment is edited.
\end{document}
